% GENERATED FILE. Edit canonical lessons, then run npm run generate:books.
% canonical-source-hash: fnv1a64:d43d69750c523192
% canonical-lessons: TE-C418-sadhyam, TE-C418-asadhyam, TE-C418-sambhavyata, TE-C418-anivaryam, TE-C418-vilaite

\chapter{Possible or Not}
\label{ch:te-possible-or-not}
\hlchaptermodality{418}

\begin{chapteropening}
\textbf{By the end of this chapter:} \emph{I can say what is possible, likely or unavoidable.}

Complete the last lesson of chapter 418: I can say what is possible, likely or unavoidable.
\end{chapteropening}

\section[\texorpdfstring{\te{సాధ్యం} (sādhyaṁ)}{సాధ్యం (sādhyaṁ)}]{\te{సాధ్యం} (sādhyaṁ) — possible; feasible}

\label{lesson:TE-C418-sadhyam}



\begin{quote}
Before the new one: say the Telugu for except for, then the Telugu for to apply.
\end{quote}

\subsection*{You'll want to know: \te{సాధ్యం}}
\textbf{\te{సాధ్యం}} — \emph{sādhyaṁ} — \textquotedblleft{}possible; feasible\textquotedblright{}.

\begin{quote}
\te{ఒక్క} \te{రోజులో} \te{ఆ} \te{పని} \te{సాధ్యం} \te{కాదు}.
\end{quote}

\begin{quote}
\emph{okka rōjulō ā pani sādhyaṁ kādu.}
\end{quote}

\begin{quote}
That job is not possible in a single day.
\end{quote}

\te{సాధ్యం} \te{కాదు}, it is not possible; \te{సాధ్యమైనంత}, as much as possible.

\begin{cousinweb}
\textbf{\te{సాధ్యం}} (\emph{sādhyaṁ}) is borrowed from Sanskrit \emph{sādhya}, \textquotedblleft{}to be mastered, won or managed\textquotedblright{} (Monier-Williams), from \emph{sādh}, \textquotedblleft{}to accomplish\textquotedblright{}.
\end{cousinweb}

\begin{grammarlens}[title={What you've built}]
One more word for saying what happens if.
\end{grammarlens}

\subsection*{Your turn}
\begin{itemize}
\raggedright
  \item \emph{Say it:} \emph{sādhyaṁ}
  \item \emph{Say it:} \emph{sādhyaṁ}, once more
  \item \emph{Say it:} \emph{okka rōjulō ā pani sādhyaṁ kādu.} — That job is not possible in a single day.
  \item \emph{Recall:} say the Telugu for except for, then the Telugu for to apply, then say \emph{sādhyaṁ} again
\end{itemize}

\begin{tcolorbox}[breakable,colback=teal!4,colframe=teal!35!black,title={Before you move on}]
What does \te{సాధ్యం} mean? (\textquotedblleft{}Possible; feasible\textquotedblright{}.) Say it once more.
\end{tcolorbox}
\section[\texorpdfstring{\te{అసాధ్యం} (asādhyaṁ)}{అసాధ్యం (asādhyaṁ)}]{\te{అసాధ్యం} (asādhyaṁ) — impossible}

\label{lesson:TE-C418-asadhyam}



\begin{quote}
Before the new one: say the Telugu for to apply, then the Telugu for possible.
\end{quote}

\subsection*{You'll want to know: \te{అసాధ్యం}}
\textbf{\te{అసాధ్యం}} — \emph{asādhyaṁ} — \textquotedblleft{}impossible\textquotedblright{}.

\begin{quote}
\te{కష్టపడితే} \te{ఏదీ} \te{అసాధ్యం} \te{కాదు}.
\end{quote}

\begin{quote}
\emph{kaṣṭapaḍitē ēdī asādhyaṁ kādu.}
\end{quote}

\begin{quote}
If you work hard, nothing is impossible.
\end{quote}

\te{అ}-, not, + \te{సాధ్యం}, possible.

\begin{cousinweb}
\textbf{\te{అసాధ్యం}} (\emph{asādhyaṁ}) is borrowed from Sanskrit \emph{asādhya}, \textquotedblleft{}not to be effected or accomplished\textquotedblright{} (Monier-Williams): \emph{a-}, \textquotedblleft{}not\textquotedblright{}, with \emph{sādhya}.
\end{cousinweb}

\begin{grammarlens}[title={What you've built}]
One more word for saying what happens if.
\end{grammarlens}

\subsection*{Your turn}
\begin{itemize}
\raggedright
  \item \emph{Say it:} \emph{asādhyaṁ}
  \item \emph{Say it:} \emph{asādhyaṁ}, once more
  \item \emph{Say it:} \emph{kaṣṭapaḍitē ēdī asādhyaṁ kādu.} — If you work hard, nothing is impossible.
  \item \emph{Recall:} say the Telugu for to apply, then the Telugu for possible, then say \emph{asādhyaṁ} again
\end{itemize}

\begin{tcolorbox}[breakable,colback=teal!4,colframe=teal!35!black,title={Before you move on}]
What does \te{అసాధ్యం} mean? (\textquotedblleft{}Impossible\textquotedblright{}.) Say it once more.
\end{tcolorbox}
\section[\texorpdfstring{\te{సంభావ్యత} (sambhāvyata)}{సంభావ్యత (sambhāvyata)}]{\te{సంభావ్యత} (sambhāvyata) — probability; likelihood}

\label{lesson:TE-C418-sambhavyata}



\begin{quote}
Before the new one: say the Telugu for possible, then the Telugu for impossible.
\end{quote}

\subsection*{You'll want to know: \te{సంభావ్యత}}
\textbf{\te{సంభావ్యత}} — \emph{sambhāvyata} — \textquotedblleft{}probability; likelihood\textquotedblright{}.

\begin{quote}
\te{రేపు} \te{వర్షం} \te{పడే} \te{సంభావ్యత} \te{ఎక్కువ}.
\end{quote}

\begin{quote}
\emph{rēpu varṣaṁ paḍē sambhāvyata ekkuva.}
\end{quote}

\begin{quote}
The likelihood of rain tomorrow is high.
\end{quote}

A careful, technical word; in talk, \te{అవకాశం}, which you know.

\begin{cousinweb}
\textbf{\te{సంభావ్యత}} (\emph{sambhāvyata}) is built on Sanskrit \emph{sambhāvya}, \textquotedblleft{}to be supposed or expected; possible, probable\textquotedblright{} (Monier-Williams), with the abstract ending -\te{త}.
\end{cousinweb}

\begin{grammarlens}[title={What you've built}]
One more word for saying what happens if.
\end{grammarlens}

\subsection*{Your turn}
\begin{itemize}
\raggedright
  \item \emph{Say it:} \emph{sambhāvyata}
  \item \emph{Say it:} \emph{sambhāvyata}, once more
  \item \emph{Say it:} \emph{rēpu varṣaṁ paḍē sambhāvyata ekkuva.} — The likelihood of rain tomorrow is high.
  \item \emph{Recall:} say the Telugu for possible, then the Telugu for impossible, then say \emph{sambhāvyata} again
\end{itemize}

\begin{tcolorbox}[breakable,colback=teal!4,colframe=teal!35!black,title={Before you move on}]
What does \te{సంభావ్యత} mean? (\textquotedblleft{}Probability; likelihood\textquotedblright{}.) Say it once more.
\end{tcolorbox}
\section[\texorpdfstring{\te{అనివార్యం} (anivāryaṁ)}{అనివార్యం (anivāryaṁ)}]{\te{అనివార్యం} (anivāryaṁ) — unavoidable}

\label{lesson:TE-C418-anivaryam}



\begin{quote}
Before the new one: say the Telugu for impossible, then the Telugu for probability.
\end{quote}

\subsection*{You'll want to know: \te{అనివార్యం}}
\textbf{\te{అనివార్యం}} — \emph{anivāryaṁ} — \textquotedblleft{}unavoidable\textquotedblright{}.

\begin{quote}
\te{ఆ} \te{మార్పు} \te{అనివార్యం}.
\end{quote}

\begin{quote}
\emph{ā mārpu anivāryaṁ.}
\end{quote}

\begin{quote}
That change is unavoidable.
\end{quote}

\te{అనివార్యంగా}, unavoidably.

\begin{cousinweb}
\textbf{\te{అనివార్యం}} (\emph{anivāryaṁ}) is borrowed from Sanskrit \emph{anivārya}, \textquotedblleft{}not to be warded off, unavoidable, irresistible\textquotedblright{} (Monier-Williams).
\end{cousinweb}

\begin{grammarlens}[title={What you've built}]
One more word for saying what happens if.
\end{grammarlens}

\subsection*{Your turn}
\begin{itemize}
\raggedright
  \item \emph{Say it:} \emph{anivāryaṁ}
  \item \emph{Say it:} \emph{anivāryaṁ}, once more
  \item \emph{Say it:} \emph{ā mārpu anivāryaṁ.} — That change is unavoidable.
  \item \emph{Recall:} say the Telugu for impossible, then the Telugu for probability, then say \emph{anivāryaṁ} again
\end{itemize}

\begin{tcolorbox}[breakable,colback=teal!4,colframe=teal!35!black,title={Before you move on}]
What does \te{అనివార్యం} mean? (\textquotedblleft{}Unavoidable\textquotedblright{}.) Say it once more.
\end{tcolorbox}
\section[\texorpdfstring{\te{వీలైతే} (vīlaitē)}{వీలైతే (vīlaitē)}]{\te{వీలైతే} (vīlaitē) — if possible}

\label{lesson:TE-C418-vilaite}



\begin{quote}
Before the new one: say the Telugu for probability, then the Telugu for unavoidable.
\end{quote}

\subsection*{You'll want to know: \te{వీలైతే}}
\textbf{\te{వీలైతే}} — \emph{vīlaitē} — \textquotedblleft{}if possible\textquotedblright{}.

\begin{quote}
\te{వీలైతే} \te{నేను} \te{రేపు} \te{మీ} \te{ఇంటికి} \te{వస్తాను}.
\end{quote}

\begin{quote}
\emph{vīlaitē nēnu rēpu mī iṇṭiki vastānu.}
\end{quote}

\begin{quote}
If possible, I will come to your house tomorrow.
\end{quote}

\te{వీలు}, a chance, + \te{అయితే}, if it is. \te{వీలు} \te{కాదు}, it is not convenient.

\begin{cousinweb}
\textbf{\te{వీలైతే}} (\emph{vīlaitē}) is built from \emph{vīlu}, \textquotedblleft{}convenience, a chance\textquotedblright{}, whose origin is uncertain, and \emph{aitē}, \textquotedblleft{}if it is\textquotedblright{}, a form of \emph{avu}, \textquotedblleft{}to be\textquotedblright{} (DEDR 333).
\end{cousinweb}

\begin{grammarlens}[title={What you've built}]
One more word for saying what happens if.
\end{grammarlens}

\subsection*{Your turn}
\begin{itemize}
\raggedright
  \item \emph{Say it:} \emph{vīlaitē}
  \item \emph{Say it:} \emph{vīlaitē}, once more
  \item \emph{Say it:} \emph{vīlaitē nēnu rēpu mī iṇṭiki vastānu.} — If possible, I will come to your house tomorrow.
  \item \emph{Recall:} say the Telugu for probability, then the Telugu for unavoidable, then say \emph{vīlaitē} again
\end{itemize}

\begin{tcolorbox}[breakable,colback=teal!4,colframe=teal!35!black,title={Before you move on}]
What does \te{వీలైతే} mean? (\textquotedblleft{}If possible\textquotedblright{}.) Say it once more.
\end{tcolorbox}
